\section{Conclusions}
\label{sec:transferability}

$\mathrm{Ba_5Y_{12}Zn[O(SiO_4)]_8}$ rests on sound single-crystal structural evidence and on evidence of synthesis by the high-temperature solution method, but no public evidence yet shows that its bulk phase-pure range has been independently reproduced~\cite{ababaikeri2024ba5y12zn}. At this stage it would be premature to prescribe a ``best recipe''; what can be specified are the variables that must be controlled together and the characterization steps that go with them.

The neighboring systems supply three consistent lessons. First, phase formation in the Ba--Y--Si--O system is sensitive to the composition coordinates, the MoO$_3$ content, and the cooling path, and the Zn-free tetragonal phase additionally involves non-stoichiometry, competing phases, and a glass-forming region~\cite{wierzbickawieczorek2017high,yamane2024synthesis,gulay2024navigation}. Second, secondary phases in ceramic samples change how the nominal composition must be interpreted, and a single-crystal assignment is no substitute for bulk phase quantification~\cite{yamane2024microstructure}. Third, the polymorphism and phase stability of the Ba--Zn--Si--O system vary with thermal history and substituted composition, so the Zn source, the cooling path, and the observation temperature must be recorded together~\cite{lin1999phase,kerstan2013bazn2si2o7,zou2019anti}. These lessons can guide experimental design, but the numerical values of the related systems cannot be carried over as conditions for the target phase.

The most informative work for the next stage is, in order: (i) establishing a local phase diagram around the target composition from BaCO$_3$, Y$_2$O$_3$, ZnO, and SiO$_2$ as the base starting materials; (ii) examining the $\mathrm{Ba_5Y_{13-\mathit{x}}Zn_{\mathit{x}}[SiO_4]_8O_{8.5-\mathit{x}/2}}$ series to distinguish limited solid solution, a line compound, and a two-phase region; (iii) running solid-state reaction and high-temperature solution growth in parallel, with the bulk phase region fixed by PXRD\slash Rietveld analysis and the structural assignment by SCXRD; and (iv) testing the size difference through the Mg\slash Co analogs only once the target phase has been independently reproduced. The oxide or carbonate combinations recommended by the retrosynthesis model can serve as a starting point for the precursors, but they cannot replace phase-diagram, thermal-history, and structural verification.

The experiment most worth starting now is a fine-step local phase diagram around the target composition: BaCO$_3$, Y$_2$O$_3$, ZnO, and SiO$_2$ as precursors, staged calcination and annealing within the temperature range of the neighboring systems, and quantitative PXRD\slash Rietveld analysis plus EDS or EPMA at every composition point, with the first aim of establishing whether the target phase has a reproducible bulk phase region; the near-single-phase compositions are then carried into high-temperature solution growth and the structure is re-checked by SCXRD. Such experiments update the hypotheses whether they succeed or fail. A fine-step composition grid, the measured element ratios and oxygen stoichiometry, cross-validation between single-crystal and powder diffraction, and fully preserved negative results together form an iterable experimental basis and bring the following questions into focus: whether Zn enters the Zn-free tetragonal silicate series, whether the target phase has a reproducible stability region, and whether the crystal-growth conditions and the bulk phase-formation conditions belong to the same phase-equilibrium region.
